\documentclass[9pt,onecolumn,twoside,lineno]{pnas-new}
% Use the lineno option to display guide line numbers if required.

% Define custom prompt environment
\definecolor{promptbg}{RGB}{250,250,250}
\definecolor{promptframe}{RGB}{200,200,200}
\newmdenv[
    linewidth=1pt,
    backgroundcolor=promptbg,
    linecolor=promptframe,
    leftmargin=1em,
    rightmargin=1em,
    innerleftmargin=1em,
    innerrightmargin=1em,
    innertopmargin=1em,
    innerbottommargin=1em,
    skipabove=1em,
    skipbelow=1em
]{prompt}


\templatetype{pnasresearcharticle} % Template for a two-column research article

\title{AI-Driven Knowledge Synthesis for Ecological Model Parameterization}

% Use letters for affiliations, numbers to show equal authorship (if applicable) and to indicate the corresponding author
\author[a,b,1]{Scott Spillias}
\author[a,b]{Elizabeth A. Fulton} 
\author[b,c,d]{Fabio Boschetti}
\author[a]{Cathy Bulman}
\author[c]{Joanna Strzelecki}
\author[a,b]{Rowan Trebilco}

\affil[a]{CSIRO Environment, Hobart, Australia}
\affil[b]{Centre for Marine Socio-Ecology, University of Tasmania, Hobart, Australia}
\affil[c]{CSIRO Environment, IOMRC Crawley, Australia}
\affil[d]{Oceans Institute, The University of Western Australia, Crawley, WA}

% Please give the surname of the lead author for the running footer
\leadauthor{Spillias} 

% Please add here a significance statement to explain the relevance of your work
\significancestatement{This study introduces an AI-driven framework for automated species grouping and food web construction for ecosystem models, addressing a critical bottleneck in model development. The framework demonstrates high reproducibility in species grouping decisions and food web construction, with strong ecological alignment when validated against expert-derived food webs. These findings demonstrate the framework's potential to generate reproducible, ecologically meaningful food webs while significantly reducing development time.}

% Please include corresponding author, author contribution and author declaration information
\authorcontributions{SS: Conceptualization, Methodology, Software, Validation, Formal analysis, Investigation, Resources, Data Curation, Writing - Original Draft, Writing - Review \& Editing, Visualization, Project administration. BF: Validation, Writing - Review \& Editing, Supervision, Funding acquisition. FB: Methodology, Software, Validation, Writing - Review \& Editing, Supervision. CB: Investigation, Data Curation, Validation, Writing - Review \& Editing. JS: Conceptualization, Validation, Investigation, Writing - Review \& Editing. RT: Methodology, Software, Validation, Investigation, Writing - Review \& Editing, Supervision, Funding acquisition.}
\authordeclaration{The authors declare no competing interest.}
\correspondingauthor{\textsuperscript{1}To whom correspondence should be addressed. E-mail: scott.spillias@csiro.au}

% Keywords are not mandatory, but authors are strongly encouraged to provide them. If provided, please include two to five keywords, separated by the pipe symbol, e.g:
\keywords{Artificial Intelligence $|$ Ecological Modelling $|$ Food Web Construction $|$ Trophic Interactions $|$ Large Language Models} 

\begin{abstract}
This study introduces a novel Large Language Model (LLM) driven framework for automated species grouping and food web construction for ecosystem models, addressing a critical bottleneck in model development. The framework (i) retrieves a marine species list from an area; (ii) uses LLMs to classify them into functional groups; and (iii) synthesises trophic interactions from diverse data sources including global biodiversity databases, species interaction repositories, and unstructured user-provided text. We evaluate the framework across four large Australian marine regions to assess both consistency and ecological accuracy of the resulting functional groups and trophic interactions. The framework demonstrates high reproducibility in species grouping decisions (>99.7\% consistency) and food web construction, with 51-59\% of predator-prey interactions showing consistent trophic proportions across multiple runs. Validation against expert-derived food webs for the Great Australian Bight ecosystem reveals strong ecological alignment and accuracy, with 92.6\% of taxonomic assignments being at least partially correct (>75\% fully correct), and correctly identifying 85\% of trophic interactions, while estimating interaction strengths within 0.2 of expert values for 80\% of interactions. These findings demonstrate the framework's potential to generate reproducible, ecologically meaningful food webs for ecosystem model development while significantly reducing development time.
\end{abstract}

\dates{This manuscript was compiled on \today}
\doi{\url{www.pnas.org/cgi/doi/10.1073/pnas.XXXXXXXXXX}}

\begin{document}

\maketitle
\thispagestyle{firststyle}
\ifthenelse{\boolean{shortarticle}}{\ifthenelse{\boolean{singlecolumn}}{\abscontentformatted}{\abscontent}}{}

% Introduction - Note: PNAS doesn't use "Introduction" as a section heading
Ecosystem modelling is an important tool for understanding and managing complex environments, with food web models being essential for representing trophic interactions and energy flows in marine ecosystems and predicting their responses to external pressures \citep{Christensen2004, Colleter2015,audzijonyte2019atlantis}. These models provide quantitative insights into ecosystem structure and function, enabling researchers to assess cumulative impacts of multiple stressors and support ecosystem-based management decisions \citep{Coll2015, Villasante2016,Geary2020}. However, constructing these models presents significant challenges, particularly in developing food webs that capture the complex web of trophic interactions within an ecosystem.

Traditional approaches to food web construction rely heavily on extensive literature review, data collation and expert knowledge, which are time-consuming and resource-intensive \citep{Holden2024a}. The process of assembling diet matrices is particularly challenging, requiring synthesis of diverse data sources including field studies, literature reviews, and expert opinion. This creates a significant bottleneck in ecosystem model development, especially when applying models to new geographical contexts \citep{Holden2024b}. Recent advances in artificial intelligence (AI) offer new opportunities to streamline the food web construction process and avoid such bottlenecks \citep{spilliasfuture}. AI tools have demonstrated success in both knowledge/evidence synthesis tasks \citep{spillias2024human,keck2025extracting,castro2024large,spillias2024evaluating,Zheng2023,nugraha2024traditional}, ecological and environmental tasks \citep{Fernandes2024,Li2024,Chen2024,dorm2025large,Noleto2024} and modelling tasks \citep{Lapeyrolerie2022, Tuia2022, Karniadakis2021}, but their application to process-based ecosystem modelling remains nascent. The key challenge lies in ensuring that AI-driven approaches can effectively synthesise available information while maintaining ecological validity.

We present a flexible framework for assembling and synthesizing user-defined and online resources to construct food webs using Large Language Models (LLMs). Our approach integrates multiple data sources, including global biodiversity databases, species interaction repositories, and locally-held unstructured or structured text, to automate key steps in food web development. The framework employs user-selected LLMs to group species into functional units and estimate trophic interaction strengths. We evaluate the system in four distinct Australian marine ecosystems - the Northern Australia, South East shelf, and South East offshore regions, where we assess the reproducibility of the approach, and in the Great Australian Bight where we assess the accuracy of the approach. Specifically, we test the precision (repeatability) and scientific accuracy of automated species grouping decisions, and the precision and accuracy of the resulting diet matrix proportions, with accuracy defined in terms of similarity to expert estimates. These regions offer contrasting environmental conditions, species assemblages, and ecological dynamics, providing a robust test of the framework's adaptability and reliability.

Our framework operates through a five-stage process: (i) species identification using the Ocean Biodiversity Information System (OBIS) to generate comprehensive species lists for defined regions; (ii) data harvesting from FishBase, SeaLifeBase, and Global Biotic Interactions (GLOBI) databases to collect ecological and trophic information; (iii) species grouping using LLMs to classify species into functional groups based on ecological roles; (iv) food web construction using retrieval-augmented generation to synthesize species-level data into group-level diet compositions; and (v) validation through multiple iterations to assess reproducibility and comparison with expert-derived models to evaluate accuracy. For reproducibility assessment, we conducted five independent iterations per region, calculating consistency scores and Jaccard similarity indices to measure grouping stability, and stability scores to evaluate trophic interaction consistency. For accuracy assessment, we compared our framework's outputs with an expert-derived Ecopath model of the Great Australian Bight ecosystem \citep{Fulton2018}, evaluating both taxonomic grouping decisions and diet matrix proportions. This methodological approach allows us to systematically evaluate both the precision and ecological validity of our AI-driven food web construction framework.




\input{results}

\input{discussion}

\matmethods{
    \input{methods}
}
\showmatmethods{} % Display the Materials and Methods section

\acknow{SS was funded by a CSIRO R+ Postdoctoral Fellowship. This project was partially funded by the `Futures of Seafood' project. Generative AI tools, specifically Claude Sonnet 3.5, were utilized in the preparation of this manuscript to assist with tasks such as language refinement, text structuring, and summarization. All scientific content, data interpretation, and conclusions were independently developed and verified by the authors to ensure accuracy and integrity.}

\showacknow{} % Display the acknowledgments section

% Bibliography
\bibliography{references}

\end{document}